\newpage
\section{Spécification du projet}
\label{sec:specification}

\subsection{Notions de base et contraintes du projet}
\label{sec:spec1}

\paragraph{Fonctionnement général du logiciel}
Le jeu se déroule en tours. Chaque joueur reçoit 5 cartes au début. 
Le joueur avec la plus petite carte (hors 2 et Joker) commence. 
À chaque tour, le joueur doit poser une combinaison qui bat la précédente. 
Si aucun joueur ne peut jouer, le tas est réinitialisé.

\paragraph{Déroulement détaillé d'une partie}
Une partie se décompose en plusieurs rounds. Chaque round permet à chaque joueur de jouer à tour de rôle. Lorsqu'un joueur ne peut pas jouer, il pioche une carte et passe son tour. Si tous les joueurs passent leur tour sans jouer (sauf le dernier ayant joué), le tas est réinitialisé : toutes les cartes sont mélangées et replacées dans la pioche. La partie se termine immédiatement lorsqu'un joueur pose sa dernière carte.

\paragraph{Notions et terminologies de base}
\begin{itemize}
\item \textbf{Carte} : 54 cartes (52 standards + 2 Jokers). Les valeurs vont du 2 à l'As.
\item \textbf{Combinaison simple} : 1 carte. C'est la combinaison la plus élémentaire.
\item \textbf{Double} : 2 cartes de valeur identique (exemple : deux Rois).
\item \textbf{Série} : 3 cartes ou plus dont les valeurs se suivent (exemple : 4,5,6).
\item \textbf{Bombe} : 3 ou 4 cartes de valeur identique. Une bombe bat toute combinaison simple.
\item \textbf{Double Joker} : 2 Jokers ensemble. C'est la combinaison la plus puissante.
\end{itemize}

\paragraph{Hiérarchie des combinaisons}
Les combinaisons sont classées par force croissante : Simple (<) Double (<) Série (<) Bombe (<) Double Joker. Une combinaison ne peut battre une autre que si elle est de force supérieure ou égale (cas des séries où les cartes doivent être plus hautes).

\paragraph{Contraintes et limitations connues}
\begin{itemize}
\item Langage : Java 8 (compatibilité ascendante)
\item Interface graphique : Swing (bibliothèque standard)
\item Bibliothèques externes : Log4j pour les logs, JUnit pour les tests, jFreeChart pour les graphiques
\item Problème connu : la musique ne fonctionne que sous Windows (bug du format WAV sous Linux)
\item L'interface est conçue pour une résolution minimale de 1200x650 pixels
\end{itemize}

\subsection{Fonctionnalités attendues du projet}
\label{sec:spec2}

\paragraph{Fonctionnalités implémentées}
L'ensemble des fonctionnalités demandées dans le cahier des charges a été réalisé :
\begin{itemize}
\item Distribution de 5 cartes à chaque joueur en début de partie
\item Détection automatique du joueur qui commence (carte la plus faible)
\item Pose de combinaisons (simple, double, série, bombe, double joker)
\item Gestion des cartes spéciales (2 écrase tout, joker polymorphe)
\item Pioche automatique si aucune carte jouable
\item Reset du tour quand tous les joueurs passent
\item Recyclage du deck quand il est vide
\item Interface graphique complète (menu, options, jeu, fin de partie)
\item 3 niveaux de difficulté pour les bots (Facile, Moyen, Difficile)
\item Système de scores (1 carte = 1 point, bombe = ×2, double joker = ×4)
\item Logs détaillés (format texte et HTML)
\item Tests unitaires JUnit (couverture des classes critiques)
\item Graphiques statistiques avec jFreeChart (camembert et barres)
\end{itemize}

\paragraph{Fonctionnalités non implémentées (hors périmètre)}
Certaines fonctionnalités n'ont pas été retenues dans le périmètre du projet :
\begin{itemize}
\item Mode multijoueur en ligne
\item Sauvegarde des parties
\item Version mobile
\end{itemize}
